% /Users/pikachu/books/payments-book/src/parts/part4/ch21-subscriptions.tex
\chapter{Подписки и Card-on-File}
\label{ch:subscriptions}
\margintoc

\begin{kaobox}[title=Что вы узнаете из этой главы]
  Подписка начинается с первого согласия на хранение реквизитов,
  а не со второго списания. Дальше важно понять, зачем схемы ввели
  фреймворк сохранённых реквизитов, как различать транзакции,
  инициированные клиентом и мерчантом, почему сетевой идентификатор
  транзакции нельзя терять при миграции, как вести себя при мягком
  отказе и чем обновление реквизитов помогает снижать технический отток.
\end{kaobox}

Представим самый будничный сценарий: в январе клиент сам оформляет
подписку, вводит карту и соглашается на регулярные списания, а в феврале
система пытается списать ту же сумму уже без его участия. Для продукта
это выглядит как естественное продолжение одной услуги, но для эмитента
это две разные авторизационные ситуации. Именно эту устойчивость
и создаёт дисциплина, с которой система связывает первое согласие
клиента со всеми последующими списаниями. Ни интерфейс автосписания,
ни календарь биллинга сами по себе эту задачу не решают. Если связь
разорвать, эмитент перестаёт понимать контекст операции, схема начинает
считать её подозрительной, а клиент сталкивается с лишними отказами там,
где платёж должен был пройти спокойно.

\section{Stored Credential Framework}
\label{sec:subscriptions-scf}

До появления единых правил мерчанты решали повторные списания каждый
по-своему: кто-то хранил PAN, кто-то полагался на токенизацию, кто-то
каждый месяц отправлял обычный запрос и надеялся, что эмитент сам
угадает контекст. Трудность была не в самом повторном платеже, а в том,
что эмитенту не сообщали главное: это продолжение уже согласованной
цепочки или новая операция без явного участия клиента.

\scheme{Visa} и \scheme{Mastercard} сформировали фреймворк сохранённых
реквизитов (Stored Credential Framework, SCF)\sidenote{Stored Credential
Framework — набор правил, который позволяет схеме различать первое
согласие клиента и последующие списания по сохранённым реквизитам.},
чтобы контекст первого согласия и последующих списаний стал видимым в
авторизационном сообщении \sidecite{visa-scf}\sidecite{mc-scf}. Идея
здесь несложная: карта остаётся у мерчанта как сохранённый реквизит, но
каждая операция должна ясно показывать, каким образом она связана с
исходным согласием клиента.

В авторизационном сообщении эта связь выражается несколькими
признаками: первичным или последующим использованием сохранённых
реквизитов, типом сценария и сетевым идентификатором транзакции,
который связывает всю цепочку. Иначе говоря, SCF нужен затем, чтобы
эмитент видел границу между первичным согласием клиента и дальнейшим
автосписанием, а не только формальный набор реквизитов.

Важно и ограничение этого подхода: SCF не делает любую повторную
операцию автоматически безопасной. Если мерчант не передаёт контекст
или путает тип операции, эмитент вправе трактовать запрос как обычную
бесконтекстную CNP-транзакцию и применить к нему более строгую политику
проверки.

\section{CIT и MIT}
\label{sec:subscriptions-cit-mit}

Подписочный процесс почти всегда состоит из двух сцен. В первой клиент
сам оформляет услугу, соглашается на сохранение реквизитов и проходит
все необходимые шаги подтверждения. Во второй мерчант через неделю или
месяц пытается списать очередной платёж уже без участия клиента. Именно
на границе между этими двумя сценами чаще всего и ломается подписочный
процесс.

Транзакция, инициированная клиентом (customer-initiated transaction,
CIT), относится к первой сцене. Клиент сам нажимает кнопку оплаты,
вводит данные карты, проходит аутентификацию, если она требуется
правилами рынка. Так выглядит первый платёж за подписку, оформление
Card-on-File или покупка, после которой мерчант получает право
сохранить реквизиты. В этот момент клиент не просто платит: система
фиксирует исходную точку цепочки и получает тот самый
сетевой идентификатор транзакции (Network Transaction ID, NTID)\sidenote{Network
Transaction ID — сетевой идентификатор, который схема возвращает в ответе
на успешную CIT и который связывает последующие MIT с исходным согласием.}.

Для CIT важно не перепутать юридический и технический смысл. Согласие
клиента действительно возникает здесь, но сохранение реквизитов само по
себе ещё не означает право на любые будущие списания. Дальше система
должна хранить NTID и передавать его в последующих операциях, иначе
связь между первым согласием и следующими списаниями быстро исчезнет из
авторизационного контекста.

Транзакция, инициированная мерчантом (merchant-initiated transaction,
MIT), относится уже ко второй сцене. Это списание без активного участия
клиента: ежемесячная подписка, доначисление по уже оказанной услуге или
платёж по ранее согласованной схеме. Здесь мерчант инициирует запрос,
но обязан показать эмитенту, что операция продолжает уже существующую
цепочку. Поэтому MIT почти всегда живёт рядом со SCF, NTID и схемными
признаками типа операции.

Ограничение здесь жёсткое: MIT без правильного контекста может быть
расценена как новая и подозрительная транзакция. В европейском трафике
это особенно заметно, потому что режимы сильной аутентификации и
исключений завязаны на корректную классификацию операции. Если
инициатор и контекст описаны неверно, эмитент вправе отклонить платёж
даже тогда, когда бизнес-логика мерчанта считает его штатным.

CIT создаёт согласие, а MIT использует уже созданную цепочку.
Для подписочного бизнеса правило простое: первый платёж
должен не только пройти, но и заложить техническую основу для всех
следующих списаний.

\section{NTID и цепочка согласия}
\label{sec:subscriptions-ntid}

Зачем вообще хранить NTID отдельно от токена или PAN? Затем, что токен
отвечает на вопрос, каким реквизитом платить, а NTID --- на вопрос,
почему это автосписание связано с уже данным согласием. Это разные
сущности, и смешивать их опасно: токен можно обновить при перевыпуске
карты, а NTID должен пережить смену реквизитов, если согласие клиента
не менялось.

Типичный сбой возникает при миграции между PSP. Мерчант переносит
токены и историю подписок, но забывает перенести NTID первой CIT. На
следующем цикле MIT приходит уже без этого якоря, и часть эмитентов
перестаёт видеть в платеже продолжение старой цепочки. В результате
падает доля одобрений, хотя с точки зрения бизнеса услуга та же самая,
а клиент тот же.

На бумаге защита выглядит просто: нужно хранить NTID как обязательный
атрибут подписки, передавать его во всех последующих MIT и не
разрывать цепочку при смене процессора. На практике это требует такой
же аккуратности, как работа с токенами и самой записью о подписке.
Если исходная CIT не была оформлена корректно, никакой последующий
косметический ремонт уже не спасёт цепочку согласия.

Отсюда практический вывод: NTID следует считать не вспомогательным
полем, а ключом к жизненному циклу подписки. Потеряли NTID --- потеряли
часть авторизационной истории.

\section{Мягкий отказ и повторная попытка}
\label{sec:subscriptions-soft-decline}

Как действовать, если MIT получает мягкий отказ? Прежде всего не
воспринимать его как случайную помеху. Мягкий отказ означает, что
эмитент не готов одобрить платёж в его текущем виде, но и не ставит на
цепочке окончательный крест. Типичный пример --- код, требующий
аутентификации: в таком случае платёж не надо гонять по кругу без
изменений.

Дальнейшая реакция зависит от причины отказа. Если эмитент просит
аутентификацию, клиенту нужен новый сеанс подтверждения, после чего
цепочка согласия может быть продолжена. Если же отказ связан не с
аутентификацией, а с временной политикой риска или нехваткой средств,
повторная попытка должна следовать правилам схемы, а не внутренней
логике самого мерчанта.

Произвольный повтор MIT вне схемного протокола только ухудшает
картину. Для эмитента такая система начинает выглядеть как источник
шума, а аккуратного биллинга из этого поведения уже не получается.
Поэтому повторная попытка допустима только тогда, когда она опирается
на понятную причину отказа и на разрешённый сценарий схемы.

Итак, мягкий отказ --- это сигнал пересмотреть контекст платежа и, если
нужно, вернуть клиента в цепочку подтверждения, а не приглашение к
бесконечному повтору.

\section{Обновление реквизитов}
\label{sec:subscriptions-account-updater}

Что делать, когда клиент не менял подписку, но карта у него уже новая?
Именно для этого нужны сервисы обновления реквизитов. Сценарий знаком:
эмитент перевыпустил карту, старый PAN больше не подходит, а мерчант
продолжает пытаться списывать по старому набору данных.

Visa Account Updater (VAU) и Mastercard Automatic Billing Updater
(ABU) позволяют получать обновлённые реквизиты от участвующих
эмитентов \sidecite{visa-core-rules}\sidecite{mc-rules}. Обычно это
работает пакетно: мерчант или его процессор передаёт список сохранённых
реквизитов, а схема возвращает актуальные данные там, где эмитент
разрешил такой обмен.

Ограничения у этого подхода тоже важны. Во-первых, эмитент может не
участвовать в сервисе обновления. Во-вторых, обновление не мгновенно:
между перевыпуском карты и получением новых реквизитов есть временная
задержка. В-третьих, если мерчант уже перешёл на сетевые токены, часть
этой работы берёт на себя сама схема, и зависимость от пакетного
обновления становится меньше.

Практический вывод здесь такой: обновление реквизитов не заменяет
дисциплину SCF и NTID, но заметно снижает технический отток и число
отказов по старым картам. Для зрелой подписочной системы это не
украшение, а базовый элемент гигиены.

\section{Как выглядит устойчивый подписочный процесс}
\label{sec:subscriptions-anti-churn}

Какая последовательность делает подписку устойчивой? Начинается всё с
первой CIT: она должна не просто успешно пройти, а корректно
зафиксировать согласие. Затем система сохраняет NTID, токен или другой идентификатор
реквизита, а клиент получает ясное описание условий подписки:
периодичности, суммы и правил изменения.

Дальше каждый биллинговый цикл строится как MIT с корректными флагами,
своевременной проверкой актуальности реквизитов и понятной обработкой
отказов. Если карта перевыпущена, помогает обновление реквизитов. Если
эмитент просит вернуть клиента в цикл подтверждения, это надо сделать.
Если отказ носит временный характер, повторная попытка должна идти по
правилам схемы. Импровизация команды здесь особенно опасна, потому что
она размывает именно тот контекст, на котором держится вся подписка.

Отдельная часть процесса --- коммуникация с клиентом. Здесь важен не
универсальный таймер, а предсказуемость: клиент должен заранее знать,
когда будет следующая попытка списания и что именно ему нужно сделать,
если платёж не прошёл. Для подписки это часто важнее, чем красивая
форма уведомления, потому что понятное ожидание снижает и отток, и
число бессмысленных повторов.

\begin{kaobox}[title=На практике]
  Устойчивость подписки обычно держится на четырёх вещах: корректная
  первая CIT, сохранённый NTID, актуальные реквизиты и схемный порядок
  обработки отказов. Если хотя бы один элемент выпадает, платёжный
  поток начинает терять авторизации, а клиент видит это как
  необъяснимый сбой сервиса.
\end{kaobox}

\begin{chaptersummary}
\begin{itemize}
  \item Stored Credential Framework сделал видимым контекст сохранённых
        реквизитов и отделил первое согласие клиента от последующих списаний.
  \item CIT создаёт исходную точку согласия, MIT продолжает её; NTID
        связывает всю цепочку и должен переживать смену процессора.
  \item Потеря NTID при миграции между PSP ведёт к падению
        доли одобрений, даже если токены и сами реквизиты сохранены.
  \item Мягкий отказ нельзя лечить произвольным повтором: сначала надо
        понять причину отказа и вернуть клиента в правильный сценарий.
  \item VAU и ABU помогают переживать перевыпуск карт, но не отменяют
        дисциплину SCF, NTID и схемной обработки отказов.
  \item Устойчивый подписочный процесс строится как цепочка, а не как
        набор разрозненных списаний.
\end{itemize}
\end{chaptersummary}
